% Prologue of The Replica Rosetta Stone (written by the coordinating author).
\specialchapter{Prologue}{The number two}{proved}
\rstfloatprefix{P}
\begin{refsection}[chapters/00-prologue.bib]

Take $n$ points at random in $d$-dimensional space and give each a random label, plus or minus. Is there a hyperplane through the origin with every plus point on one side and every minus point on the other? If $n$ is much smaller than $d$ the answer is surely yes, and if $n$ is much larger it is surely no. What is surprising is how sharp the boundary is, and where it lies. When $d$ is large, the answer switches from almost certainly yes to almost certainly no as the ratio $\alpha=n/d$ crosses $2$.

The number $2$ has been derived at least four times, in four languages. A combinatorialist counts the labelings that a hyperplane can realize. A physicist computes the volume of the set of separating hyperplanes with replicas and a saddle point. A probabilist compares a Gaussian process with a simpler one. A geometer measures the size of a cone. The four derivations agree, and they reach their common answer by such different roads that a reader fluent in one can hardly follow another. This book is about numbers like this one: exact predictions about inference and learning in high dimension, derived first, in most cases, by statistical physicists, and proved later, in many cases, by mathematicians. It sets the derivations side by side, step by step, and records what has been proved, what is believed, and what is known to be false.

\section*{Counting}
\addcontentsline{toc}{section}{Counting}

Write $z_\mu=y_\mu x_\mu$ for the $\mu$-th point multiplied by its label. The labels are separable exactly when some $w\in\R^d$ has $\ip{z_\mu}{w}>0$ for every $\mu\le n$. Because the points are Gaussian and the labels are random signs, the $z_\mu$ are themselves independent standard Gaussian vectors, so the question is whether $n$ random Gaussian vectors lie in a common open half-space through the origin.

\textcite{Cover1965} answered it by counting. Call $n$ points in $\R^d$ in general position if every $d$ of them are linearly independent, and let $C(n,d)$ be the number of the $2^n$ labelings of the points that a hyperplane through the origin can separate. Cover showed that $C(n,d)$ depends on the points only through $n$ and $d$, and that it satisfies the recursion $C(n+1,d)=C(n,d)+C(n,d-1)$: a separable labeling of the first $n$ points extends to the new point in both ways exactly when some separating hyperplane passes through the new point, which is a separation problem in one dimension fewer. With $C(1,d)=2$ the recursion gives
\begin{equation}\label{eq:prologue:cover}
\begin{gathered}
C(n,d)=2\sum_{k=0}^{d-1}\binom{n-1}{k},\\
\Prob(\text{separable})=\frac{C(n,d)}{2^n}=\Prob\bigl(\mathrm{Bin}(n-1,\tfrac12)\le d-1\bigr).
\end{gathered}
\end{equation}
The count of the regions that $n$ hyperplanes through the origin cut out of $\R^d$ goes back to Schläfli, and the probability in \eqref{eq:prologue:cover}, for every symmetric distribution in general position, was computed by \textcite{Wendel1962}. At $n=2d$ the binomial is symmetric about $d-\frac12$, so the probability is exactly $\frac12$; for $n=\alpha d$ with $d\to\infty$ it tends to $1$ when $\alpha<2$ and to $0$ when $\alpha>2$, across a window of width of order $\sqrt d$. \Cref{fig:prologue:cover} shows the formula at four sizes against linear programs on simulated data.

\begin{figure}[t]
\centering
\begin{tikzpicture}
\begin{axis}[rplot,xlabel={$\alpha=n/d$},ylabel={probability of separability},xmin=1,xmax=3,ymin=-0.03,ymax=1.03,
  legend pos=north east,legend style={font=\scriptsize}]
\addplot[proof!35,line width=0.8pt,no marks] table[x=alpha,y=d25] {figures/data/00-prologue-cover-curves.dat};
\addplot[proof!55,line width=0.8pt,no marks] table[x=alpha,y=d50] {figures/data/00-prologue-cover-curves.dat};
\addplot[proof!80,line width=0.9pt,no marks] table[x=alpha,y=d100] {figures/data/00-prologue-cover-curves.dat};
\addplot[proof,line width=1.0pt,no marks] table[x=alpha,y=d200] {figures/data/00-prologue-cover-curves.dat};
\addplot[measured,error bars/.cd,y dir=both,y explicit] table[x=alpha,y=frac,y error plus=hi,y error minus=lo] {figures/data/00-prologue-cover-sim.dat};
\addplot[guide] coordinates {(2,-0.03) (2,1.03)};
\legend{$d=25$,$d=50$,$d=100$,$d=200$,linear programs ($d=100$)}
\end{axis}
\end{tikzpicture}
\caption{Cover's exact probability \eqref{eq:prologue:cover} that $n=\alpha d$ random labeled Gaussian points can be separated by a hyperplane through the origin, at $d=25$, $50$, $100$ and $200$ (lapis curves, darker for larger $d$), and the fraction of $60$ simulated instances at $d=100$ that a linear program finds separable, with $95\%$ Clopper--Pearson intervals (seed $2026$). Every interval contains the exact value. The curves cross $\frac12$ at $\alpha=2$ for every $d$ and sharpen as $d$ grows. Script: \texttt{figures/scripts/00-prologue-cover.py}.}\label{fig:prologue:cover}
\end{figure}

Counting is exact at every size, which no other derivation in this book will be. It is also fragile: ask for a margin, so that every point must lie at distance at least $\kappa$ from the hyperplane, or restrict the weights to $\pm1$, and the count no longer depends only on $n$ and $d$.

\section*{Replicas}
\addcontentsline{toc}{section}{Replicas}

\textcite{Gardner1988} asked the question for every margin at once, and in the language of statistical physics. Normalize the weights to the sphere $\norm w^2=d$ and the patterns to $g_\mu\sim\Normal(0,I_d)$, and ask for the fraction of the sphere on which all $n$ constraints $\ip{g_\mu}{w}/\sqrt d\ge\kappa$ hold. This volume $V$ is a random variable that depends on the patterns, and its logarithm, divided by $d$, is what a physicist calls an entropy. It should concentrate around its mean, and the capacity $\alpha_c(\kappa)$ is the load at which the mean entropy falls to $-\infty$: the solutions disappear.

\begin{physicsbox}[Gardner 1988]
The capacity of the spherical perceptron with margin $\kappa$ is
\[
\alpha_c(\kappa)=\Bigl[\int_{-\kappa}^{\infty}(z+\kappa)^2\Dz\Bigr]^{-1}=\frac1{(1+\kappa^2)\Phi(\kappa)+\kappa\phi(\kappa)},
\qquad\alpha_c(0)=2 .
\]
\end{physicsbox}

\begin{heuristic}
The mean of $\log V$ is hard, but the mean of $V^r$ is not, for integer $r$: it is the volume of $r$ copies, or \emph{replicas}, $w^1,\dots,w^r$ that all satisfy the same constraints. For fixed replicas, the fields $\ip{g_\mu}{w^a}/\sqrt d$ of one pattern are jointly Gaussian with covariance given by the overlaps $Q_{ab}=\ip{w^a}{w^b}/d$, and the $n$ patterns are independent, so $\E V^r$ is an integral over the overlap matrix of $\exp(d\,S(Q))$ for an explicit function $S$. For large $d$ the integral is dominated by a saddle point. Guess that it is \emph{replica symmetric}, $Q_{ab}=q$ for $a\ne b$; continue the result from integer $r$ to $r\to0$, and use $\E\log V=\lim_{r\to0}(\E V^r-1)/r$. The typical entropy becomes
\[
\frac1d\E\log V=\min_{0\le q<1}\Bigl[\alpha\,\E\log H\Bigl(\frac{\kappa+\sqrt q\,Z}{\sqrt{1-q}}\Bigr)+\frac q{2(1-q)}+\frac12\log(1-q)\Bigr],
\]
with $H(x)=\Prob(Z>x)$; here $q$ is the overlap of two typical solutions. As $\alpha$ grows the solutions shrink toward a point and $q\to1$. Since $\log H(x)\approx-x^2/2$ for large $x$, the bracket behaves as $(1-\alpha\E[(Z+\kappa)_+^2])/(2(1-q))$, which stays finite only while $\alpha\E[(Z+\kappa)_+^2]\le1$. Equality is the capacity.
\end{heuristic}

Every step after the first is unjustified by the standards of a mathematician: the continuation from integer $r$ to $r\to0$, the exchange of $r\to0$ with $d\to\infty$, a saddle point that is a minimum over $q$ where Laplace's method would ask for a maximum, and above all the guess that the saddle point is replica symmetric. Yet the answer is right for every $\kappa\ge0$. For $\kappa<0$ the same computation shows its own failure: the replica-symmetric saddle point becomes unstable at the capacity, and physicists predict that the true answer requires breaking the replica symmetry completely \parencite{FranzParisi2016,FranzEtAl2017}.

\section*{Comparison}
\addcontentsline{toc}{section}{Comparison}

The probabilist's route starts from the observation that feasibility is the sign of a Gaussian min--max. With unit vectors $u=w/\sqrt d$ and the matrix $G$ whose rows are the patterns, the constraints hold for some $u$ exactly when
\[
\xi=\min_{\norm u\le1}\ \max_{\lambda\ge0,\ \norm\lambda\le1}\ \lambda^\top(\kappa\mathbf1-Gu)
\]
is at most $0$, at least when $\kappa>0$, because then a solution can be rescaled onto the sphere. The function inside is bilinear in the Gaussian matrix $G$. \textcite{Gordon1985} proved that the minimum of the maximum of such a process can be compared, in probability, with that of a much simpler one, in which $\lambda^\top Gu$ is replaced by $\norm u\,g^\top\lambda+\norm\lambda\,h^\top u$ with independent standard Gaussian vectors $g\in\R^n$ and $h\in\R^d$. The simple problem can be solved by hand: its value changes sign when $\norm{(\kappa\mathbf1+g)_+}=\norm h$, that is, by the law of large numbers, when $n\,\E[(Z+\kappa)_+^2]=d$. This is Gardner's formula. For $\kappa\ge0$ every set in sight is convex, and the comparison holds in both directions (\cref{ch:gordon}); the formula is then a theorem, proved by \textcite{Stojnic2013a} in this way, and earlier by \textcite{ShcherbinaTirozzi2003} and by \textcite[Chapter~8]{Talagrand2011b} through the volume itself. For $\kappa<0$ only one direction survives: the formula is an upper bound on the capacity, and \textcite{Stojnic2013b} showed that it is not the capacity for $\kappa\lesssim-0.63$.

\section*{Geometry}
\addcontentsline{toc}{section}{Geometry}

The fourth route is the shortest. The labels are separable when the $d$-dimensional subspace spanned by the columns of the $n\times d$ matrix $Z$ meets the nonnegative orthant of $\R^n$ in more than the origin, and the subspace is uniformly distributed because $Z$ is Gaussian. \textcite{AmelunxenEtAl2014} proved that a random subspace of dimension $d$ meets a convex cone $C\subset\R^n$ with high probability if $d+\delta(C)>n+O(\sqrt n)$ and misses it if $d+\delta(C)<n-O(\sqrt n)$, where the statistical dimension $\delta(C)=\E\norm{\Pi_C(g)}^2$ is the mean squared length of the projection of a Gaussian vector onto the cone. The orthant is its own polar cone, and for such a cone a Gaussian vector splits into two projections of equal mean squared length, so $\delta=n/2$. The threshold $d+n/2=n$ is $n=2d$.

\section*{One number, two languages}
\addcontentsline{toc}{section}{One number, two languages}

Of the four derivations, two will recur throughout the book: the physicist's replicas and the mathematician's comparison. Set side by side, they are not as foreign to each other as they look.

\begin{stone}{Gardner's replicas}{Gordon's comparison}
\stonestep{Set up}{The typical volume of the solutions, $\frac1d\E\log V$, through the moments $\E V^r$ of $r$ replicas.}{The sign of a Gaussian min--max $\xi$, bilinear in the patterns.}
\stonestep{Reduce}{The replicas interact only through their overlaps $Q_{ab}$; replica symmetry leaves one number $q$.}{Gordon's comparison replaces the Gaussian matrix by two Gaussian vectors; the problem reduces to two norms, of $u$ and of $\lambda$.}
\stonestep{Solve}{At capacity $q\to1$, and the divergent terms balance when $\alpha\,\E[(Z+\kappa)_+^2]=1$.}{The auxiliary problem changes sign when $\norm{(\kappa\mathbf1+g)_+}=\norm h$, that is, when $\alpha\,\E[(Z+\kappa)_+^2]=1$.}
\stonestep{Check}{The replica-symmetric saddle point is stable for $\kappa\ge0$ and unstable for $\kappa<0$, where replica symmetry breaks.}{Convexity gives the comparison in both directions for $\kappa\ge0$; for $\kappa<0$ only Gordon's one-sided bound remains.}
\end{stone}

The rows match. The physicist's single overlap $q$ and the probabilist's two norms carry the same information, the same equation appears at the end of both derivations, and the place where the physicist distrusts his own answer, $\kappa<0$, is exactly the place where the probabilist's tool stops working. \Cref{ch:gordon} shows that this is no accident: for convex problems, the replica-symmetric saddle point and Gordon's auxiliary problem are the same variational problem, which is why the physicists were right there, and why the comparison theorem says nothing more where they had doubts. \Cref{fig:prologue:margin} measures the largest margin of random labels against Gardner's formula.

\begin{figure}[t]
\centering
\begin{tikzpicture}
\begin{axis}[rplot,xlabel={$\alpha=n/d$},ylabel={largest margin $\kappa$},xmin=0,xmax=4,ymin=-1.2,ymax=2.1,legend pos=north east]
\addplot[physpred] table[x=alpha,y=kappa] {figures/data/00-prologue-margin-curve.dat};
\addplot[physpred2] table[x=alpha,y=kappa] {figures/data/00-prologue-margin-curveneg.dat};
\addplot[measured,ebars] table[x=alpha,y=kappa,y error plus=half,y error minus=half] {figures/data/00-prologue-margin.dat};
\addplot[guide] coordinates {(0,0) (4,0)};
\legend{Gardner's formula (proved),Gardner's formula (an upper bound only),largest margin ($d=200$)}
\end{axis}
\end{tikzpicture}
\caption{The largest margin $\kappa_{\max}=\max_{\norm u=1}\min_\mu\ip{z_\mu}{u}$ of $n=\alpha d$ random labeled Gaussian points in dimension $d=200$, computed exactly as the distance from the origin to their convex hull ($20$ instances per point, Student-$t$ $95\%$ intervals, seed $7$), against the inverse of Gardner's formula, $\alpha_c(\kappa)=\alpha$. For $\alpha<2$ the formula is a theorem; six of the seven intervals contain it, and the seventh, at $\alpha=1.5$, misses by $0.008$, within the finite-size fluctuations at $d=200$. For $\alpha>2$ the largest margin is negative, the problem is not convex, and the dotted curve is only an upper bound. Script: \texttt{figures/scripts/00-prologue-cover.py}.}\label{fig:prologue:margin}
\end{figure}

\section*{The status board}
\addcontentsline{toc}{section}{The status board}

Every chapter of this book ends its work on a prediction with a board that says, in a word for each statement, how much of it is known. The words are \emph{proved}, for a theorem in a refereed publication; \emph{proved under} extra hypotheses, which the board names; \emph{preprint}, for a proof that so far exists only in an unrefereed preprint; \emph{conjectured}, with the evidence; \emph{heuristic}, for an argument that claims no exactness; \emph{known false}, with the counterexample; and \emph{open}. The perceptron's board, as of October 2026, reads as follows.

\begin{statusboard}
\statusrow{Capacity $2$ at margin $0$, at every finite size}{\proved}{\textcite{Cover1965}; \textcite{Wendel1962}}
\statusrow{Gardner's formula $\alpha_c(\kappa)$ for $\kappa\ge0$}{\proved}{\textcite{ShcherbinaTirozzi2003}; \textcite{Stojnic2013a}; \textcite[Chapter~8]{Talagrand2011b}}
\statusrow{Gardner's formula for $\kappa<0$, as an upper bound}{\proved}{\textcite{Stojnic2013a}}
\statusrow{Gardner's formula for $\kappa<0$, as the capacity}{\knownfalse}{for $\kappa\lesssim-0.63$ \parencite{Stojnic2013b}; physics predicts full replica-symmetry breaking for every $\kappa<0$ \parencite{FranzEtAl2017}}
\statusrow{The capacity for $\kappa<0$}{\openq}{known to leading order as $\kappa\to-\infty$ \parencite{MontanariZhongZhou2024}}
\statusrow{Weights in $\{\pm1\}^d$: capacity $0.83308$}{\provedunder}{numerical conditions \parencite{DingSun2019,Huang2024}; conjectured by \textcite{KrauthMezard1989}}
\statusrow{The same, without conditions}{\provedpreprint}{computer-assisted proof, arXiv:2607.11010v1, which its author says has not yet been independently reproduced or peer reviewed \parencite{Shmalo2026}}
\end{statusboard}

\section*{What this book does}
\addcontentsline{toc}{section}{What this book does}

The perceptron is the first of fifteen predictions. \emph{Part I} introduces the two languages on the simplest disordered systems. \Cref{ch:freeentropy} treats free entropies, saddle points and self-averaging, with the Gibbs measure read as a Bayesian posterior; \cref{ch:rem} the random energy model, where the replica trick can be checked against an exact computation; and \cref{ch:sk} the Sherrington--Kirkpatrick model, where replica symmetry breaks and Guerra's interpolation explains why the physicists' answer is nevertheless right. \emph{Part II} turns to Bayesian inference, where the posterior is exactly the Gibbs measure that physics studies and replica symmetry becomes a theorem: the Nishimori line (\cref{ch:nishimori}), spiked matrices and the BBP transition (\cref{ch:spiked}), message passing (\cref{ch:amp}), and generalized linear models with their hard phase (\cref{ch:glm}). \emph{Part III} is about convex learning problems, where Gordon's comparison replaces the replicas: Gaussian comparison and the perceptron (\cref{ch:gordon}), ridge regression (\cref{ch:ridge}), the LASSO and robust regression (\cref{ch:lasso}), logistic regression and max-margin classifiers (\cref{ch:classification}), and kernels with the Gaussian equivalence principle (\cref{ch:kernels}). \emph{Part IV} turns from what can be learned to where algorithms go: the critical points of random landscapes (\cref{ch:kacrice}), dynamical mean-field theory (\cref{ch:dmft}), and online stochastic gradient descent (\cref{ch:sgd}).

Every chapter has the same anatomy. \emph{The question} states the prediction in a carnelian box. \emph{The physics} derives it as a physicist would, honestly, naming every uncontrolled step. \emph{The stone} sets the derivation and the rigorous route side by side, row by row, across the full width of the page, as above. \emph{The proof} proves the central result completely, in general or in a special case that keeps the idea, and states the general theorems with their sources. \emph{Measuring it} runs a short script, available with the book, and compares the measurement with the prediction (carnelian) and with what is proved (lapis). \emph{The status board} lists every prediction of the chapter with its status, and \emph{Where the translation fails} says what was assumed, where the physics is wrong or not yet right, and what is open. Exercises and notes close the chapter; the notes translate the notation of the main sources of both literatures into the book's.

\paragraph{How to read.} The book can be read in order, but it is written for two kinds of readers. A reader trained in probability may want to read each \emph{The physics} section slowly, as a new language, and to use the stone as a dictionary. A reader trained in physics may want to start from the stone and read \emph{The proof} as an account of what a mathematician must add. A short course is \cref{ch:freeentropy,ch:sk,ch:nishimori,ch:amp,ch:gordon,ch:ridge}, followed by the chapters of Part~IV, each of which can be read on its own. The prerequisites are linear algebra, multivariable calculus, probability at the level of a first graduate course, and some convex optimization; the appendix (\cref{ch:toolkit}) states and proves the Gaussian tools and states the random-matrix facts that the book uses.

\paragraph{What this book does not do.} It does not prove the Parisi formula, whose proof fills books of its own \parencite{Talagrand2011b,Panchenko2013b}; it states it precisely and proves the parts that the rest of the book uses. It treats mean-field models, in which every variable interacts with every other, and not models on finite-dimensional lattices. It does not claim to be the first text to put replica computations next to proofs: \textcite{MontanariSen2024} and the lecture notes of \textcite{KrzakalaZdeborova2024} do so for selected models, and \textcite{BandeiraPerryWein2018} explain the physicists' heuristics to mathematicians. To our knowledge it is the first systematic parallel text of this kind for inference and learning that includes Gaussian comparison, landscapes and dynamics, with a status board for every prediction.

\section*{Notes}
\addcontentsline{toc}{section}{Notes}

The function counting theorem is \textcite{Cover1965}; \textcite{Wendel1962} computed the probability that symmetric random points lie in a common half-space. The capacity of the perceptron with a margin is due to \textcite{Gardner1988}, whose space-of-interactions method founded the statistical mechanics of learning; \textcite{GardnerDerrida1988} extended it to other costs and to binary weights, where the replica-symmetric answer is wrong and \textcite{KrauthMezard1989} found the correct value. Gordon's comparison inequality is \textcite{Gordon1985}; its use for exact asymptotics, with convex duality, is due to \textcite{Stojnic2013a}, and the two-sided convex form, the convex Gaussian min--max theorem, to \textcite{ThrampoulidisOymakHassibi2015} and \textcite{ThrampoulidisAbbasiHassibi2018}. The kinematic formula for random cones and the statistical dimension are from \textcite{AmelunxenEtAl2014}. For the physics of the perceptron below zero margin, see \textcite{FranzParisi2016} and \textcite{FranzEtAl2017}; for rigorous bounds there, \textcite{MontanariZhongZhou2024}. The textbooks of the two literatures are \textcite{MezardParisiVirasoro1987}, \textcite{Nishimori2001}, \textcite{EngelVandenBroeck2001} and \textcite{MezardMontanari2009} on one side, and \textcite{Talagrand2011a,Talagrand2011b} and \textcite{Panchenko2013b} on the other; the review of \textcite{ZdeborovaKrzakala2016} surveys the statistical physics of inference.

\printchapterbib
\end{refsection}
\rstfloatrestore
